\section{Experiments}\label{sec:experiments}

We evaluated the three envelope construction methods on simulated data with missing scores and on two football 
classification tasks. We keep $\alpha=0.1$ throughout, the sample of each class is split equally into the shape discovery 
set $S_1$ and the size scaling set $S_2$, finally coverage is the proportion of test samples whose prediction set contains 
the true label. We report coverage, the average set size, the singleton rate and the empty-set rate, as mean $\pm$ 
standard error over repeated runs. The radial method uses $M=100$, $\kappa=4$ and $\delta=30^\circ$, and the strip 
method uses $NB=8$ with $n_0=3$ on the simulated data and $n_0=10$ on the football data; all other parameters are the 
default parameters in the library. The code and the saved outputs are available at \citep{conformalenvelopes}.
% \url{https://github.com/Nizalia13/Conformal_Classification_Library}.

\paragraph{Simulated data.} There are two classes, $A$ and $B$, and four predictors which are used to score each candidate 
class. For a sample with true class $h^\ast$ and a candidate class $h$, let $\sigma=+1$ if $h=h^\ast$ and $\sigma=-1$ 
otherwise. The first predictor has the score $1/(1+e^{-(1.8\,\sigma+\varepsilon_1)})$ with $\varepsilon_1\sim N(0,1)$. 
The second is built in the same way from a noise term $\varepsilon_2$ with correlation $0.9$ to $\varepsilon_1$. The 
third is weak, it has a shift $0.35\,\sigma$ and noise of standard deviation $1.5$, and the fourth is uniform on $[0,1]$, 
it carries no information. Each score is then independently hidden with probability $0.15$. If all four scores of a 
candidate are hidden the draw is repeated. In each run there are $1500$ training draws and $500$ test samples, and we use 
$500$ independent runs.

\paragraph{Football data.} We classify FIFA~19 players from $34$ skill ratings \citep{fifa19data}. In the first task 
the label is simpler, it is just the role (defender, midfielder, attacker or goalkeeper) for $18{,}147$ players. In the 
second it is the recorded position, in this we only use the $17{,}835$ players in the $20$ positions with at least $100$ 
player records. Four classifiers (logistic regression, random forest, gradient boosting and a support vector classifier) 
are trained on a separate $40\%$ of the players, and the scores are their predicted probabilities of each label. Another
$20\%$ is used for validation, and the remaining players are used for conformal fitting ($30\%$) and testing ($10\%$). 
To measure the variability of the coverage, we keep the trained classifiers fixed and evaluate them over $100$ random split 
of the remaining data into calibration and test sets (sized $5{,}444$ / $1{,}815$ for roles and 
$5{,}351$ / $1{,}784$ for positions). The data contain no missing scores.

 
\begin{table}[ht]
\centering
\small
\begin{tabular}{llcccc}
\toprule
Data & Method & Coverage & Set size & Singleton & Empty \\
\midrule
Simulated & Collapsed & $0.9012\pm0.0008$ & $1.014$ & $0.811$ & $0.088$ \\
($500$ runs) & Radial & $0.9001\pm0.0008$ & $1.086$ & $0.752$ & $0.081$ \\
 & Strip & $0.8998\pm0.0008$ & $0.955$ & $0.856$ & $0.095$ \\
\midrule
Football roles & Collapsed & $0.9002\pm0.0009$ & $1.069$ & $0.887$ & $0.022$ \\
($100$ runs) & Radial & $0.8995\pm0.0010$ & $1.068$ & $0.886$ & $0.023$ \\
 & Strip & $0.8994\pm0.0010$ & $1.082$ & $0.856$ & $0.031$ \\
\midrule
Football positions & Collapsed & $0.9032\pm0.0009$ & $4.364$ & $0.108$ & $0.010$ \\
($100$ runs) & Radial & $0.9016\pm0.0009$ & $4.738$ & $0.101$ & $0.010$ \\
 & Strip & $0.9016\pm0.0009$ & $4.575$ & $0.118$ & $0.011$ \\
\bottomrule
\end{tabular}
\caption{Results as mean $\pm$ standard error over repeated runs, with target coverage $0.9$.}
\label{tab:results}
\end{table}

\paragraph{Results.} Table~\ref{tab:results} shows that the coverage of all three constructions is close to the target 
$0.9$ in every experiment. The coverage of each class is also close to the target; it is between $0.899$ and $0.902$  on 
the simulated data, and between $0.896$ and $0.902$  for the four roles. For the positions, the average coverage of a 
position lies between $0.888$ and $0.926$, it has more range due to the smaller number of data points for each position.

The methods have similar set sizes on the football role tasks, where most of the sets are singletons. But on the position 
classification task the sets contain between $4.4$ and $4.7$ positions on average, which is expected, because neighbouring 
positions, such as the left, central and right variants of a role, have very similar skill values. The order of the methods 
by set size changes between the experiments, therefore we cannot claim that one construction is better in general. The 
football dataset is complete, so the handling of missing scores is only evaluated on the simulated data.